\subsection{SUBGROUPS}

DEFINITION 4. Let G be a group with operators. A stable subgroup of G is a subset H of G with the following properties:

\begin{enumerate}
\def\labelenumi{(\roman{enumi})}
\item
  \(e\in \mathbf{H};\)
\item
  \(x, y \in \mathbf{H}\) implies \(xy \in \mathbf{H}\) ;
\item
  \(x \in \mathbf{H}\) implies \(x^{-1} \in \mathbf{H}\) ;
\item
  \(x \in \mathbf{H}\) and \(\alpha \in \Omega\) imply \(x^{\alpha} \in \mathbf{H}\) .
\end{enumerate}

If H is a stable subgroup of G, the structure induced on H by the structure of a group with operators on G is the structure of a group with operators and the canonical injection of H into G is a homomorphism of groups with operators.

Let G be a group. A stable subgroup of G with the action of \(\varnothing\) (no. 2), which is a subset of G satisfying conditions (i), (ii), (iii) of Definition 4, is called a subgroup of G. When we speak of a subgroup of a group of operators we shall always mean a subgroup of the underlying group of G. A subgroup of a group with operators G is not necessarily a stable subgroup of G.

Example (1). Let \(\Sigma\) be a species of structure (Set Theory, IV, § 1, no. 4) and S a structure of species \(\Sigma\) on a set E (loc. cit.). The set of automorphisms of S is a subgroup of \(\mathfrak{S}_{\mathbb{E}}\) .

PROPOSITION 1. Let G be a group with operators and H a subset of G which is stable under the homotheties of G. The following conditions are equivalent:

\begin{enumerate}
\def\labelenumi{(\alph{enumi})}
\item
  H is a stable subgroup of G.
\item
  H is non-empty and the relations \(x \in \mathbf{H}, y \in \mathbf{H}\) imply \(xy \in \mathbf{H}\) and \(x^{-1} \in \mathbf{H}\) .
\item
  H is non-empty and the relations \(x \in \mathbf{H}\) , \(y \in \mathbf{H}\) imply \(xy^{-1} \in \mathbf{H}\) .
\item
  H is stable under the law on G and the law of composition induced on H by the law of composition on G is a group law.
\end{enumerate}

Clearly (a) implies (b). We show that (b) implies (a). It suffices to show that H contains the identity element of G. As the subset H is non-empty, let \(x \in H\) . Then \(x^{-1} \in H\) and \(e = xx^{-1} \in H\) . Clearly (b) implies (c). We show that (c) implies (b). First of all since H is non-empty it contains an element x. Hence \(xx^{-1} = e\) is an element of H. For every element \(x\) of H, \(x^{-1} = ex^{-1}\) belongs to H; hence the relations \(x \in H\) , \(y \in H\) imply \(x(y^{-1})^{-1} = xy \in H\) . Clearly (a) implies (d). We show that (d) implies (a): the canonical injection of H into G is a group homomorphism; hence \(e \in H\) and the relation \(x \in H\) implies \(x^{-1} \in H\) (no. 1).

Remarks. (1) Similarly it can be shown that condition (b) is equivalent to the condition

(c') H ≠ ∅ and the relations \(x \in H\) and \(y \in H\) imply \(y^{-1}x \in H\) .

\begin{enumerate}
\def\labelenumi{(\arabic{enumi})}
\setcounter{enumi}{1}
\tightlist
\item
  For every subgroup H of G there are the following relations
\end{enumerate}

\[
\mathrm{H.H} = \mathrm{H} \quad \text { and } \quad \mathrm{H} ^ {- 1} = \mathrm{H}.\tag{1}
\]

For H.H ⊂ H and H \(^{-1}\) ⊂ H by (b). As e \ensuremath{\in} H, H.H ⊃ e.H = H and taking inverses transforms the inclusion H \(^{-1}\) ⊂ H into H ⊂ H \(^{-1}\) , whence formulae (1).

If H is a stable subgroup of G and K is a stable subgroup of H, clearly K is a stable subgroup of G.

The set \(\{e\}\) is the smallest stable subgroup of G. The intersection of a family of stable subgroups of G is a stable subgroup. There is therefore a smallest stable subgroup H of G containing a given subset X of G; it is called the stable subgroup generated by X and X is called a generating system (or generating set) of H.

PROPOSITION 2. Let X be a non-empty subset of a group with operators G and \(\hat{X}\) the stable subset under the action of \(\Omega\) on G generated by X. The stable subgroup generated by X is the stable subset under the law on G generated by the set Y = \(\hat{X} \cup \hat{X}^{-1}\) .

The latter subset Z is the set of compositions of finite sequences all of whose terms are elements of \(\hat{X}\) or inverse of elements of \(\hat{X}\) : the inverse of such a composition is a composition of the same form (§ 2, no. 3, Corollary 1 to Proposition 5) and Z is stable under the action of \(\Omega\) , as is seen by applying § 3, no. 4, Proposition 1 to the homotheties of G, hence (Proposition 1) Z is a stable subgroup of G. Conversely, every stable subgroup containing X obviously contains Y and hence Z.

COROLLARY 1. Let G be a group with operators and X a subset of G which is stable under the action of Ω. The subgroup generated by X and the stable subgroup generated by X coincide.

COROLLARY 2. Let G be a group and X a subset of G consisting of pairwise permutable elements. The subgroup of G generated by X is commutative.

The set \(Y = X \cup X^{-1}\) consists of pairwise permutable elements (§2, no. 3, Proposition 5) and the law induced on the stable subset generated by Y is commutative (§1, no. 5, Corollary 2).

If G is a group with operators, the stable subgroup generated by a subset of G consisting of pairwise permutable elements is not necessarily commutative.

COROLLARY 3. Let \(f: \mathbf{G} \to \mathbf{G}'\) be a homomorphism of groups with operators and \(\mathbf{X}\) a subset of \(\mathbf{G}\) . The image under \(f\) of the stable subgroup of \(\mathbf{G}\) generated by \(\mathbf{X}\) is the stable subgroup of \(\mathbf{G}'\) generated by \(f(\mathbf{X})\) .

Let \(\mathbf{X}' = f(\mathbf{X})\) . Then \(\hat{\mathbf{X}}' = f(\hat{\mathbf{X}})\) and \(\mathbf{X}^{\prime -1} = f(\mathbf{X}^{-1})\) . Hence

\[
f (\hat {\mathbf {X}} \cup \hat {\mathbf {X}} ^ {- 1}) = \hat {\mathbf {X}} ^ {\prime} \cup \hat {\mathbf {X}} ^ {- 1}.
\]

The corollary then follows from § 1, no. 4, Proposition 1.

Example (2). Let G be a group and x an element of G. The subgroup generated by \(\{x\}\) (called more simply the subgroup generated by x) is the set of \(x^{n}, n \in Z\) . The stable subset (under the law on G) generated by \(\{x\}\) is the set of \(x^{n}\) where \(n \in N^{*}\) . These two sets are in general distinct.

Thus, in the additive group \(\mathbf{Z}\) , the subgroup generated by an element \(x\) is the set \(x.\mathbf{Z}\) of \(xn\) , \(n \in \mathbf{Z}\) , and the stable subset generated by \(x\) is the set of \(xn\) , \(n \in \mathbf{N}^*\) . These two sets are always distinct if \(x \neq 0\) .

The union of a right directed family of stable subgroups of G is obviously a stable subgroup. It follows that, if P is a subset of G and H a stable subgroup of G not meeting P, the set of stable subgroups of G containing H and not meeting P, ordered by inclusion, is inductive (Set Theory, III, § 2, no. 4). Applying Zorn's Lemma (Set Theory, III, § 2, no. 4), we obtain the following result:

PROPOSITION 3. Let G be a group with operators, P a subset of G and H a stable subgroup G not meeting P. The set of stable subgroups of G containing H and not meeting P has a maximal element.

